\chapter{АНАЛИЗ ЗАРУБЕЖНЫХ И ОТЕЧЕСТВЕННЫХ ИСТОЧНИКОВ ПО ТЕМЕ ИССЛЕДОВАНИЯ} \label{ch:анализ_зарубежных_и_отечественных_источников_по_теме_исследования}

\section{Анализ отечественных источнико} \label{sec:анализ_отечественных_источников}

Исследование Е.~Н.~Платонова и И.~Р.~Мартыновой~\cite{Platonov2024} посвящено семантическому анализу отзывов об организациях с применением методов машинного обучения на русскоязычном корпусе, насчитывающем свыше 500,000 текстов. Сравнительному тестированию подвергается широкий спектр архитектур: от классических алгоритмов TF–IDF совместно с логистической регрессией, CatBoost, SVM и мультиномиальным наивным байесом — до нейросетевых подходов на основе CNN, LSTM и трансформеров. Наилучший результат по взвешенной F1–мере зафиксирован у архитектуры rubert–tiny2: значение метрики составило~0,83 при времени обучения 60~минут. Сокращение числа слоёв энкодера до одного удерживает F1–меру на отметке~0,81, уменьшая время обучения до 10~минут — тем самым подтверждается значительный потенциал облегчённых трансформерных архитектур в ресурсоограниченных средах. Показательна нестабильность полнослойной модели ruBERT Base: в идентичных условиях обучения она достигает F1–меры лишь~0,26 вследствие переобучения. Алгоритм LIME, применённый для интерпретации результатов rubert–tiny2, обеспечивает визуализацию вкладов отдельных лексем в итоговое предсказание — что существенно для понимания логики классификатора. Полученные в работе выводы непосредственно применимы при выборе базовой архитектуры для дообучения, поскольку подтверждают преимущество компактных трансформеров перед крупными моделями в условиях ограниченных вычислительных ресурсов.

В исследовании Т.~М.~Татарниковой и Н.~С.~Мокрецова~\cite{Tatarnikova2025} рассматривается подход к дистилляции языковых моделей, при котором модель-учитель подключается к процессу обучения избирательно, а не на постоянной основе. Ключевое отличие предложенного метода SwitchLLM от классических SGO-схем (Student–Generated Outputs) состоит в механизме переключения: переход от модели ученика к учителю инициируется лишь тогда, когда их вероятностные распределения расходятся сверх допустимого предела, определяемого дивергенцией Дженсена–Шеннона с экспоненциально снижающимся порогом. Сравнительное тестирование на корпусах Dolly и S–NI показало, что SwitchLLM удерживает метрику Rouge–L на длинных последовательностях практически без потерь, тогда как у MiniLLM и DistiLLM аналогичный показатель падает примерно вдвое. Примечательно и то, что среднеквадратичная ошибка (MSE) трёхмиллиардной модели-ученика выходит на уровень, сопоставимый с показателем семимиллиардной модели-учителя. Изложенные в работе теоретические принципы дистилляции непосредственно применимы к компактным архитектурам класса DistilBERT, задействованным в экспериментальной части настоящего исследования.

П.~В.~Стрюкова, Д.~В.~Герберта и П.~В.~Кожевниковой~\cite{Stryukov2025} проводят систематический анализ техник квантования LLM, разграничивая два принципиально разных сценария: PTQ, применяемое без изменения весов после обучения, и QAT, встраивающее квантование непосредственно в тренировочный процесс. В рамках PTQ авторы подробно разбирают алгоритмы GPTQ, AWQ и SpQR, каждый из которых по-своему учитывает статистику весов и активаций при их сжатии. С практической точки зрения замена формата FP16 на INT8 даёт двух-трёхкратный прирост скорости матричных операций там, где аппаратура поддерживает целочисленную арифметику, тогда как дальнейшее снижение разрядности до 2–3~бит заметно ухудшает перплексию модели. QAT демонстрирует существенное преимущество перед PTQ по сохранению точности — потери сокращаются на 50–70\%, — однако расплачивается за это ростом вычислительной нагрузки при обучении на 30–50\% и более. Особого рассмотрения удостаивается метод QLoRA: объединяя 4–битное представление весов в формате NF4 с адаптерами LoRA, он делает возможным дообучение 65-миллиардной модели на единственном GPU с 48~ГБ видеопамяти. Эксперименты на семействе LLaMA–3 демонстрируют закономерность: компактные модели теряют качество при квантовании интенсивнее, чем крупные, что объясняется большей концентрацией знаний на один параметр. Изложенные данные обосновывают выбор квантованных моделей для браузерного развёртывания: именно форматы INT8 и INT4 лежат в основе аппаратного ускорения ONNX Runtime Web, применяемого в данной работе.

З.~М.~Акиева, Д.~А.~Мурзин и С.~И.~Эльтаев~\cite{Akieva2024} исследуют применение WebAssembly для повышения производительности интерактивных веб–приложений. Эксперименты на задачах поиска простых чисел алгоритмом «Решето Эратосфена», манипулирования большими массивами данных и рендеринга 3D–графики показывают: WASM обеспечивает трёхкратное ускорение вычислений относительно аналогичного кода на JavaScript, сокращение потребления памяти на~36\% и прирост частоты кадров на~66\%. Авторы описывают принципы работы WebAssembly — компиляцию исходного кода на языках C, C++, Rust или AssemblyScript в бинарный формат .wasm, его загрузку и трансляцию браузером в машинный код, а также взаимодействие с JavaScript через API. Поддержка инструкций SIMD и механизм Web Workers для параллельных вычислений анализируются отдельно. Среди инструментальных средств рассматриваются Emscripten, AssemblyScript и Rust; заключительный раздел посвящён перспективам интеграции WebAssembly с WebGPU. Работа формирует теоретическое обоснование применения WebAssembly как вычислительного бэкенда ONNX Runtime Web для нейросетевого вывода в браузере — центрального элемента архитектуры разрабатываемой системы.

З.~А.~Овчаров~\cite{Ovcharov2023} рассматривает методы оптимизации браузерных вычислений на основе распараллеливания потоков, сопоставляя стандартный подход на базе JavaScript, механизм Web Workers и систему WebAssembly применительно к задаче поиска пересечений на множестве HTML–элементов. Оба альтернативных механизма линейно превосходят стандартный подход по времени выполнения и времени блокировки основного потока; при этом WASM демонстрирует наивысшую производительность. Детально анализируется однопоточная асинхронная модель выполнения JavaScript и механизм Event Loop: показано, что очередь функций отложенного вызова не устраняет блокировок при наличии высоконагруженных участков кода. Реализация задачи на Rust с последующей компиляцией через wasm–pack и вызовом модуля из JavaScript посредством динамического импорта рассматривается как практический пример организации асинхронного WASM–вывода. Работа обосновывает архитектурное решение об изоляции нейросетевого вывода в отдельном потоке Web Worker во избежание блокировок пользовательского интерфейса.

\section{Анализ зарубежных источников} \label{sec:анализ_зарубежных_источников}

F.~Faisal и A.~Anastasopoulos~\cite{faisal2024crosslingual} обратились к проблеме кросс–лингвального переноса в мультиязычных моделях, разработав стратификационную методику оценки того, насколько эффективно модель обобщает знания на новые языки без каких-либо дополнительных этапов тонкой настройки. Центральным сценарием в этой работе выступает zero–shot transfer — режим, особенно востребованный для языков, по которым размеченные обучающие данные практически недоступны. Эмпирические результаты указывают на то, что как выбор архитектуры, так и стратегия предобучения заметно влияют на итоговое качество межъязыкового обобщения. Данное наблюдение непосредственно учитывалось при выборе базовой модели в настоящей работе, ориентированной на мультиязычный анализ тональности.
 
J.~Wu и соавторы~\cite{wu2024prealign} ввели метод PreAlign, в котором согласование многоязычных представлений закладывается не на этапе дообучения, а непосредственно в ходе предобучения на разноязычных корпусах. Как свидетельствуют полученные авторами результаты, столь ранняя интеграция выравнивания обеспечивает более надёжные межъязыковые соответствия по сравнению со сценарием, где аналогичная процедура откладывается до стадии fine-tuning. Применительно к адаптации модели cardiffnlp/twitter--xlm--roberta--base--sentiment на языках с развитой морфологией этот вывод обретает конкретный практический смысл: качество классификации тональности оказывается в прямой зависимости от того, каким образом были инициализированы веса на самом раннем этапе тренировки.

Z.~Yao и соавторы~\cite{yao2022zeroquant} предложили метод ZeroQuant~– подход к квантованию трансформерных моделей без дополнительного обучения (post–training quantization), обеспечивающий существенное сокращение вычислительных затрат при минимальных потерях точности. Пословное квантование весовых матриц с точностью INT8, как показано авторами, ускоряет вывод крупных моделей, сохраняя значения метрики качества на уровне, сопоставимом с исходными FP32–результатами. Применительно к задаче браузерного вывода приведённые данные подкрепляют целесообразность восьмибитного квантования при конвертации моделей в формат ONNX.

Z.~Yao и соавторы~\cite{yao2022zeroquant} разработали ZeroQuant – технику PTQ для трансформерных архитектур, позволяющую сжимать модели без каких-либо дополнительных тренировочных итераций. Центральная идея метода состоит в пословном INT8-квантовании весовых матриц: такая гранулярность, по данным авторов, даёт ощутимый выигрыш в скорости инференса крупных моделей, удерживая метрики качества вблизи исходных FP32-значений. Для целей настоящей работы, где модели исполняются непосредственно в браузере, этот результат служит обоснованием выбора восьмибитного квантования при экспорте в формат ONNX.
 
X.~Wu и соавторы~\cite{wu2023int4} анализируют особенности квантования в формате INT4, рассматривая задержки вывода, возможность композиции с другими методами оптимизации и типичные сценарии деградации точности. Четырёхбитное квантование, согласно полученным результатам, достигает наибольшего ускорения на операциях матричного умножения, однако на ряде задач классификации сопровождается значимым снижением точности. Указанная закономерность непосредственно соотносится с задачей настоящего исследования по оценке влияния квантования на точность мультиязычного анализа тональности в браузерной среде.
 
A.~Bhandare и соавторы~\cite{bhandare2019int8} рассматривают квантование трансформерных моделей машинного перевода до формата INT8 с учётом специфики аппаратных ускорителей на базе инструкций VNNI. Описывается методика калибровки диапазонов активаций; при этом устанавливается, что корректно откалиброванное восьмибитное квантование обеспечивает ускорение вывода без значимой потери качества перевода. Полученные результаты подтверждают применимость аналогичной процедуры калибровки при подготовке ONNX–моделей для ONNX Runtime Web.
 
J.~Li и соавторы~\cite{li2023fp8bert} исследуют квантование модели BERT до формата FP8 как компромисс между полной точностью FP16/FP32 и целочисленными форматами INT4/INT8. Применение FP8, по данным авторов, сохраняет числовую устойчивость операций внимания при одновременном сокращении объёма памяти, занимаемой весовыми параметрами. Приведённые данные расширяют аналитическую базу при сравнении форматов квантования для мультиязычных моделей Xenova/bert–base–multilingual–uncased–sentiment.
 
T.~Udagawa и соавторы~\cite{udagawa2023distillation} проводят сравнительный анализ методов дистилляции знаний, не зависящих от конкретной задачи (task–agnostic), применительно к компрессии трансформерных языковых моделей; рассматриваются подходы, использующие дистилляцию на уровне промежуточных представлений, выходных логитов и весов внимания, — при этом устанавливается, что совместная дистилляция нескольких уровней обеспечивает наилучшее соотношение точности и размера модели. Данные результаты непосредственно характеризуют природу Xenova/distilbert–base–multilingual–cased–sentiments–student как студенческой модели, получаемой из мультиязычного BERT посредством многоуровневой дистилляции.
 
Y.~Huang и соавторы~\cite{huang2024kdnlp} посвятили исследование применению дистилляции знаний из крупных языковых моделей для повышения точности моделей меньшего размера на задачах NLP. Предлагаются стратегии обучения, при которых студенческая модель воспроизводит не только выходные распределения, но и промежуточные признаковые представления учительской модели. Применение таких стратегий при дообучении компактных трансформеров, ориентированных на браузерный вывод, существенно сокращает разрыв в точности по отношению к полноразмерным аналогам.
 
T.~Zhang и R.~Brannvall~\cite{zhang2025inhibidistilbert} представили модель InhibiDistilBERT, в архитектуре которой классические операции умножения внутри механизма внимания замещены операциями сложения и ReLU–активации. За счёт этого вычислительная нагрузка при выводе снижается при сохранении конкурентоспособной точности — обстоятельство, приобретающее особое значение в среде WebAssembly, где стоимость операций с плавающей точкой существенно выше, чем на серверных GPU.

C.~F.~Ruan и соавторы~\cite{ruan2024webllm} представили WebLLM~– систему высокопроизводительного вывода больших языковых моделей непосредственно в браузере с использованием WebGPU и компилятора Apache TVM. Надлежащая организация конвейера компиляции под конкретную WebGPU–конфигурацию, как показано авторами, выводит скорость генерации на уровень, сопоставимый с нативными реализациями на потребительском оборудовании. Из этого следует принципиальная осуществимость полноценного нейросетевого вывода на стороне клиента, что служит основанием для выбора ONNX Runtime Web с бэкэндом WebGPU применительно к задачам анализа тональности.
 
Q.~Wang и соавторы~\cite{wang2024anatomizing} провели детальный анализ производительности вывода моделей глубокого обучения в веб–браузерах, исследуя узкие места WebGL–, WebGPU– и WebAssembly–бэкэндов. WebGPU, согласно полученным данным, обеспечивает наибольшее ускорение на матричных операциях; WebAssembly с инструкциями SIMD, напротив, демонстрирует более стабильную производительность на устройствах без поддержки современных GPU–стандартов. Выявленное разграничение определяет критерии выбора бэкэнда вывода в зависимости от целевой аппаратной конфигурации — обстоятельство, непосредственно значимое для экспериментальной валидации настоящей работы.
 
F.~Jia и соавторы~\cite{jia2024empowering} предложили подход JIT–оптимизации ядер нейронных сетей для браузерного вывода на граничных устройствах. Суть метода состоит в динамической компиляции вычислительных ядер под конкретную топологию модели и аппаратные возможности целевого устройства, что обеспечивает существенное сокращение задержек относительно статически скомпилированных реализаций. При запуске мультиязычных трансформеров в среде Web Workers данный подход приобретает дополнительную актуальность: накладные расходы на межпотоковую передачу данных суммируются с задержками собственно вывода модели, и снижение каждой из составляющих оказывается критически важным.
 
J.~Xu и соавторы~\cite{xu2024ondevice} представили обзор языковых моделей, ориентированных на исполнение непосредственно на устройстве пользователя; охватываются методы компрессии~– квантование, дистилляция, структурное прореживание~– а также архитектурные решения, снижающие требования к памяти и вычислительным ресурсам. Браузерные развёртывания выделяются в отдельную категорию, специфику которой определяют ограниченный объём WebAssembly–памяти и отсутствие низкоуровневого доступа к GPU. Проведённый обзор закладывает теоретическое основание для обоснования отбора компактных мультиязычных моделей в настоящей работе.
 
S.~Mohammad и соавторы~\cite{mohammad2024edgeflex} описывают EdgeFlex--Transformer~--- среду выполнения трансформерных моделей, спроектированную с расчётом на граничные устройства, где вычислительные ресурсы жёстко ограничены. Отличительная черта системы~--- адаптивное управление глубиной вычислительного графа: в зависимости от текущей нагрузки и заданного бюджета задержки среда на лету перестраивает исполнение модели, балансируя между скоростью отклика и точностью предсказаний. По своей постановке эта работа перекликается с экспериментальной частью настоящего исследования, где измеряется влияние степени квантования на латентность инференса в браузерном окружении.

L.~Guo и соавторы~\cite{guo2023sti} представляют систему STI, в которой инференс NLP–моделей на граничных устройствах организован по конвейерному принципу: слои модели гибко распределяются между несколькими потоками выполнения в зависимости от текущей загрузки. Авторы показывают, что грамотное разбиение вычислительного графа на параллельно исполняемые сегменты даёт существенный выигрыш в латентности при работе с последовательностями произвольной длины. Этот вывод напрямую связан с архитектурным решением, принятым в данной работе,~--- построением многопоточного конвейера инференса на основе Web Workers.
 
Z.~Fu и соавторы~\cite{fu2023energynlp} исследуют методы адаптации задач NLP для энергоэффективного вывода на граничных устройствах с гетерогенными архитектурами памяти, предлагая стратегию размещения весовых параметров в разных уровнях памяти с целью снижения энергопотребления на мобильных и встроенных платформах. Полученные данные дополняют экспериментальную базу настоящей работы при оценке потребления ресурсов в условиях браузерного исполнения модели.
 
M.~Krasitskii и соавторы~\cite{krasitskii2025sentiment} проводят сравнительную оценку методов анализа тональности на корпусах основных европейских языков и арабского языка, где классические подходы на основе SVM и словарных методов сопоставляются с современными трансформерными моделями. Устанавливается, что морфологическая сложность языка существенно влияет на эффективность мультиязычных моделей при классификации тональности. Данный вывод непосредственно соотносится с задачей кросс–языкового тестирования на корпусах русского, английского и арабского языков в рамках настоящего исследования.
 
T.~Filip и соавторы~\cite{filip2024finetuning} рассматривают дообучение мультиязычных моделей для анализа тональности публикаций в социальных сетях на восточноевропейских языках группы V4. Дообучение на небольших размеченных корпусах целевых языков, согласно полученным данным, существенно повышает точность по сравнению с нулевым переносом; при этом наибольший прирост фиксируется для агглютинативных языков с развитой словоизменительной системой — морфологическая насыщенность которых создаёт повышенную нагрузку на механизм нулевого переноса. Из этого вытекает обоснование необходимости дообучения базовой архитектуры distilbert–base–multilingual–cased на размеченных данных с целью повышения точности на русском языке.
 
Md.~A.~Hasan~\cite{hasan2024ensemble} предложил ансамблевый подход к мультиязычному анализу тональности, при котором совместное использование нескольких трансформерных моделей повышает точность за счёт усреднения вероятностных распределений предсказаний по сравнению с отдельными моделями. Объединение моделей XLM–RoBERTa и mBERT, как показано, обеспечивает наилучшие результаты на многоязычных бенчмарках. Применение ансамблирования в браузерной среде сопряжено с дополнительными вычислительными издержками, поэтому результаты данной работы служат теоретическим обоснованием для сравнительного анализа единственной модели и ансамбля в задачах настоящего исследования.
 
A.~Y.~Radwan и соавторы~\cite{radwan2024tinymlnlp} представили схему семантической беспроводной классификации тональности на основе TinyML–подхода с сохранением приватности пользовательских данных, при которой лёгкая архитектура обеспечивает анализ тональности непосредственно на устройстве без передачи текста на внешний сервер при объёме параметров, допустимом для ресурсно–ограниченных сред. Модель приватности, реализуемая TinyML–подходом, концептуально совпадает с архитектурным принципом настоящей работы: весь конвейер вывода выполняется локально в браузере пользователя без обращения к удалённым API.